% ============================================================
\chapter{Introduction to Markov Chains}
\label{ch:markov}
% ============================================================

\section{Introduction}

Markov chains\index{Markov chain} are one of the most fundamental stochastic models.
They describe the random evolution of a system whose future depends on the past only
through the present (the Markov property). Their applications span physics (random
walks), biology (population dynamics), computer science (Monte Carlo algorithms,
PageRank), and finance.

\begin{intuition}[title=\textbf{The Markov property: memorylessness}]
Imagine a token on a board that moves randomly. At each step, the probability of
moving to a neighboring cell depends only on the current position, not on the past
trajectory. This is the Markov property: ``the future is independent of the past,
given the present.''
\end{intuition}

% ------------------------------------------------------------
\section{Definitions}
% ------------------------------------------------------------

\begin{definition}[Markov chain]
\label{def:markov_chain}
\index{Markov chain}
Let $E$ be a countable set (the \emph{state space}). A stochastic process
$(X_n)_{n \geq 0}$ taking values in $E$ is a \emph{Markov chain} if for all
$n \geq 0$ and all states $i_0, \ldots, i_{n-1}, i, j \in E$:
\[
    \PP(X_{n+1} = j \mid X_n = i, X_{n-1} = i_{n-1}, \ldots, X_0 = i_0)
    = \PP(X_{n+1} = j \mid X_n = i),
\]
whenever the conditioning is well-defined ($\PP(X_n = i, \ldots) > 0$).
\end{definition}

\begin{definition}[Homogeneous chain and transition matrix]
\label{def:transition_matrix}
\index{transition matrix}
The chain is \emph{homogeneous} (or \emph{time-homogeneous}) if the transition
probabilities do not depend on $n$:
\[
    p_{ij} = \PP(X_{n+1} = j \mid X_n = i), \qquad \forall n \geq 0.
\]
The matrix $P = (p_{ij})_{i,j \in E}$ is called the \emph{transition matrix}. It
satisfies:
\begin{enumerate}
    \item[(i)] $p_{ij} \geq 0$ for all $i, j$;
    \item[(ii)] $\sum_{j \in E} p_{ij} = 1$ for all $i$ (stochastic matrix).
\end{enumerate}
\end{definition}

\begin{definition}[Initial distribution]
The \emph{initial distribution} is the law of $X_0$, i.e.\ the vector
$\mu_0 = (\PP(X_0 = i))_{i \in E}$. The law of the chain is entirely determined
by $\mu_0$ and $P$.
\end{definition}

\begin{proposition}[Joint distribution]
For all $i_0, i_1, \ldots, i_n \in E$:
\[
    \PP(X_0 = i_0, X_1 = i_1, \ldots, X_n = i_n)
    = \mu_0(i_0)\, p_{i_0 i_1}\, p_{i_1 i_2} \cdots p_{i_{n-1} i_n}.
\]
\end{proposition}

\begin{example}[Random walk on $\Z$]
\index{random walk}
Let $X_0 = 0$ and $X_{n+1} = X_n + \xi_{n+1}$ where the $(\xi_n)$ are i.i.d.\ with
$\PP(\xi = +1) = p$ and $\PP(\xi = -1) = 1 - p$. This is a Markov chain on $E = \Z$
with $p_{i,i+1} = p$ and $p_{i,i-1} = 1 - p$.
\end{example}

\begin{example}[Ehrenfest model]
$2N$ particles are distributed between two boxes. At each step, a particle is chosen
uniformly at random and moved to the other box. If $X_n$ is the number of particles
in the first box, $(X_n)$ is a Markov chain on $\{0, 1, \ldots, 2N\}$ with
$p_{k,k-1} = k/(2N)$ and $p_{k,k+1} = 1 - k/(2N)$.
\end{example}

% ------------------------------------------------------------
\section{$n$-step transition probabilities}
% ------------------------------------------------------------

\begin{definition}
The \emph{$n$-step transition probabilities} are:
\[
    p_{ij}^{(n)} = \PP(X_n = j \mid X_0 = i).
\]
\end{definition}

\begin{theorem}[Chapman-Kolmogorov equation]
\label{thm:chapman_kolmogorov}
\index{Chapman-Kolmogorov equation}
For all $m, n \geq 0$ and all $i, j \in E$:
\[
    p_{ij}^{(m+n)} = \sum_{k \in E} p_{ik}^{(m)}\, p_{kj}^{(n)}.
\]
In matrix notation: $P^{(m+n)} = P^m \cdot P^n$, i.e.\ $P^{(n)} = P^n$.
\end{theorem}

\begin{proof}
By the law of total probability:
\begin{align*}
    p_{ij}^{(m+n)} &= \PP(X_{m+n} = j \mid X_0 = i)
    = \sum_{k \in E} \PP(X_{m+n} = j \mid X_m = k)\, \PP(X_m = k \mid X_0 = i) \\
    &= \sum_{k \in E} p_{kj}^{(n)}\, p_{ik}^{(m)}.
\end{align*}
\end{proof}

% ------------------------------------------------------------
\section{Classification of states}
% ------------------------------------------------------------

\begin{definition}[Accessibility and communication]
\index{communication}
\begin{itemize}
    \item State $j$ is \emph{accessible} from $i$, written $i \to j$, if there
    exists $n \geq 0$ such that $p_{ij}^{(n)} > 0$.
    \item States $i$ and $j$ \emph{communicate}, written $i \leftrightarrow j$, if
    $i \to j$ and $j \to i$.
\end{itemize}
The relation $\leftrightarrow$ is an equivalence relation. Its classes are called
\emph{communicating classes}.
\end{definition}

\begin{definition}[Irreducible chain]
\index{irreducible}
The chain is \emph{irreducible} if all states communicate, i.e.\ there is a single
communicating class.
\end{definition}

\begin{definition}[Return time and recurrence]
\index{recurrence}\index{transience}
The \emph{first return time} to $i$ is $T_i = \inf\{n \geq 1 : X_n = i\}$.
\begin{itemize}
    \item State $i$ is \emph{recurrent} if $\PP_i(T_i < \infty) = 1$.
    \item State $i$ is \emph{transient} if $\PP_i(T_i < \infty) < 1$.
\end{itemize}
\end{definition}

\begin{theorem}[Characterization of recurrence]
\label{thm:recurrence}
State $i$ is recurrent if and only if $\sum_{n=0}^{\infty} p_{ii}^{(n)} = +\infty$.
State $i$ is transient if and only if $\sum_{n=0}^{\infty} p_{ii}^{(n)} < +\infty$.
\end{theorem}

\begin{proof}
Let $N_i = \sum_{n=1}^{\infty} \ind_{\{X_n = i\}}$ be the number of returns to $i$.
Then $\E_i[N_i] = \sum_{n=1}^{\infty} p_{ii}^{(n)}$. Moreover,
$\PP_i(N_i \geq k) = \PP_i(T_i < \infty)^k$, so $N_i$ has a geometric distribution.
If $i$ is recurrent, $\PP_i(T_i < \infty) = 1$, so $\E_i[N_i] = +\infty$.
If $i$ is transient, $\PP_i(T_i < \infty) < 1$, so
$\E_i[N_i] = \PP_i(T_i < \infty)/(1 - \PP_i(T_i < \infty)) < +\infty$.
\end{proof}

\begin{proposition}[Class property]
Recurrence and transience are class properties: if $i$ communicates with $j$, then
$i$ is recurrent if and only if $j$ is recurrent.
\end{proposition}

\begin{definition}[Absorbing state]
\index{absorbing state}
A state $i$ is \emph{absorbing} if $p_{ii} = 1$. An absorbing state is recurrent.
\end{definition}

\begin{definition}[Period]
\index{period}
The \emph{period} of state $i$ is $d(i) = \gcd\{n \geq 1 : p_{ii}^{(n)} > 0\}$.
State $i$ is \emph{aperiodic} if $d(i) = 1$. The period is a class property.
\end{definition}

% ------------------------------------------------------------
\section{Stationary distributions}
% ------------------------------------------------------------

\begin{definition}[Stationary distribution]
\label{def:stationary}
\index{stationary distribution}
A probability vector $\pi = (\pi_i)_{i \in E}$ (with $\pi_i \geq 0$ and
$\sum_i \pi_i = 1$) is a \emph{stationary distribution} (or \emph{invariant
distribution}) for $P$ if:
\[
    \pi P = \pi, \qquad \text{i.e.} \quad
    \pi_j = \sum_{i \in E} \pi_i\, p_{ij}, \quad \forall j \in E.
\]
\end{definition}

\begin{theorem}[Existence for irreducible positive recurrent chains]
\label{thm:existence_stationary}
Let $(X_n)$ be an irreducible recurrent Markov chain. Then:
\begin{enumerate}
    \item There exists an invariant measure $\pi$ (unique up to a multiplicative
    constant) given by $\pi_j = 1/\E_j[T_j]$.
    \item $\pi$ is a probability distribution (i.e.\ $\sum_j \pi_j = 1$) if and only
    if the chain is \emph{positive recurrent}, i.e.\ $\E_i[T_i] < \infty$ for every
    (equivalently, for some) state $i$.
\end{enumerate}
\end{theorem}

\begin{remark}
If the state space $E$ is finite and the chain is irreducible, then it is
automatically positive recurrent and admits a unique stationary distribution.
\end{remark}

% ------------------------------------------------------------
\section{Ergodic theorem}
% ------------------------------------------------------------

\begin{theorem}[Ergodic theorem for Markov chains]
\label{thm:ergodic}
\index{ergodic theorem}
Let $(X_n)$ be an irreducible, positive recurrent, aperiodic Markov chain with
stationary distribution $\pi$. Then for every initial state $i$ and every state $j$:
\[
    \lim_{n \to \infty} p_{ij}^{(n)} = \pi_j.
\]
Moreover, for every bounded function $f : E \to \R$:
\[
    \frac{1}{n} \sum_{k=0}^{n-1} f(X_k)
    \xrightarrow[n \to \infty]{\text{a.s.}} \sum_{j \in E} f(j)\, \pi_j.
\]
\end{theorem}

\begin{keyformulas}[title=\textbf{Key Markov chain formulas}]
\begin{itemize}
    \item \textbf{Chapman-Kolmogorov:} $P^{(m+n)} = P^m \cdot P^n$
    \item \textbf{Distribution at time $n$:} $\mu_n = \mu_0 P^n$
    \item \textbf{Stationarity:} $\pi P = \pi$
    \item \textbf{Link with return time:} $\pi_j = 1/\E_j[T_j]$
    \item \textbf{Convergence:} $\lim_{n \to \infty} P^n = \mathbf{1} \pi^T$
    (if irreducible, positive recurrent, aperiodic)
\end{itemize}
\end{keyformulas}

% ------------------------------------------------------------
\section{Reversibility and detailed balance}
% ------------------------------------------------------------

\begin{definition}[Reversibility]
\index{reversibility}\index{detailed balance}
A Markov chain with transition matrix $P$ and stationary distribution $\pi$ is
\emph{reversible} if the \emph{detailed balance equations} are satisfied:
\[
    \pi_i\, p_{ij} = \pi_j\, p_{ji}, \qquad \forall i, j \in E.
\]
\end{definition}

\begin{proposition}
If $\pi$ satisfies the detailed balance equations, then $\pi$ is a stationary
distribution.
\end{proposition}

\begin{proof}
Summing over $i$:
$\sum_i \pi_i\, p_{ij} = \sum_i \pi_j\, p_{ji} = \pi_j \sum_i p_{ji} = \pi_j$,
since $\sum_i p_{ji} = 1$.
\end{proof}

\begin{example}[Random walk on a graph]
Let $G = (V, E)$ be a connected finite undirected graph. The simple random walk on $G$
has transition matrix $p_{ij} = 1/\mathrm{deg}(i)$ if $(i,j)$ is an edge. The
stationary distribution is $\pi_i = \mathrm{deg}(i)/(2\abs{E})$, and the detailed
balance equations are satisfied.
\end{example}

\begin{attention}[title=\textbf{Detailed balance $\neq$ stationarity}]
Detailed balance is a \emph{sufficient} but not necessary condition for stationarity.
There exist Markov chains with a stationary distribution that does not satisfy detailed
balance (e.g.\ chains with a directed cycle).
\end{attention}

% ------------------------------------------------------------
\section{Exercises}
% ------------------------------------------------------------

\begin{exercise}
Let $(X_n)$ be the symmetric random walk on $\Z$ ($p = 1/2$). Show that every state
is recurrent.
\emph{Hint}: use $p_{00}^{(2n)} = \binom{2n}{n} 2^{-2n} \sim (\pi n)^{-1/2}$.
\end{exercise}

\begin{exercise}
Show that the symmetric random walk on $\Z^3$ is transient.
\emph{Hint}: show that $\sum_n p_{00}^{(n)} < \infty$.
\end{exercise}

\begin{exercise}
Let $E = \{1, 2, 3\}$ and
$P = \begin{pmatrix} 0 & 1/2 & 1/2 \\ 1/3 & 1/3 & 1/3 \\ 1/2 & 0 & 1/2 \end{pmatrix}$.
\begin{enumerate}
    \item Verify that the chain is irreducible and aperiodic.
    \item Compute the stationary distribution $\pi$.
    \item Is the chain reversible?
\end{enumerate}
\end{exercise}

\begin{exercise}[Queueing model]
\index{queue}
Let $N \geq 1$ be the capacity of a queue. At each step, a customer arrives with
probability $\lambda$ and a customer is served with probability $\mu$ (if the queue
is non-empty), independently. Let $X_n$ be the number of customers at time $n$.
Show that $(X_n)$ is a Markov chain, write its transition matrix, and find its
stationary distribution when $\lambda < \mu$.
\end{exercise}

\begin{exercise}[Metropolis-Hastings algorithm]
\index{Metropolis-Hastings}
Let $\pi$ be a target probability distribution on a finite set $E$ and let $Q$ be an
irreducible transition matrix (the \emph{proposal}). Define the transition matrix:
\[
    p_{ij} = q_{ij} \min\!\left(1, \frac{\pi_j q_{ji}}{\pi_i q_{ij}}\right)
    \quad (i \neq j), \qquad
    p_{ii} = 1 - \sum_{j \neq i} p_{ij}.
\]
Show that $\pi$ is a stationary distribution for $P$ and that detailed balance holds.
\end{exercise}

\begin{exercise}
Let $(X_n)$ be an irreducible Markov chain on a finite state space $E$ with
$\abs{E} = N$. Show that there exists a unique stationary vector $\pi$ and that
$\pi_i > 0$ for all $i$.
\end{exercise}
